\section{Main Lean Theorem}
\label{sec:main-theorem}

The preferred public theorem is located in
\[
\code{Stokes/Global/CoverIndexedZeroCompactRepresentedStokesFirstPrinciples.lean}.
\]
It is re-exported through the slim public cover-indexed import.  This section
records the theorem shape and the internal route.

\subsection{Ambient Lean context}

The theorem is stated for a space \(M\) charted over a model \(H\), with model
with corners
\[
  I : \code{ModelWithCorners Real (Fin (n + 1) -> Real) H}.
\]
The surrounding Lean context includes the usual manifold and topological
assumptions used by the route, expressed using the typeclass and structure
mechanisms of Lean and mathlib
\cite{demoura2015lean,demoura2021lean4,baanen2022instance,mathlib2020}:
\[
\code{IsManifold I top M},\quad
\code{T2Space M},\quad
\code{SigmaCompactSpace M},
\]
finite dimensionality of the coordinate model, and the half-space model
assumption
\[
\code{HalfSpaceModel I}.
\]
The half-space model is a background geometric assumption: it tells Lean that
the model range is the upper half-space and supports the automatic construction
of pointwise boundary and interior chart boxes.

\subsection{First-principles input}

The public input structure is:
\begin{lstlisting}[language=Lean]
structure CoverIndexedZeroCompactRepresentedStokesFirstPrinciplesInput
    (I : ModelWithCorners Real (Fin (n + 1) -> Real) H)
    (omega : ManifoldForm I M n) where
  chartwiseSmooth : ManifoldForm.ChartwiseSmooth I omega
  compactSupport : IsCompact (closure (ManifoldForm.support I omega))
\end{lstlisting}
There is no public selected cover, no partition field, no boundary ambient
field, no collar-prism certificate, no local Stokes field, and no endpoint
adapter.  These objects are generated internally.

The compact carrier is:
\begin{lstlisting}[language=Lean]
def supportCarrier : Set M :=
  closure (ManifoldForm.support I omega)
\end{lstlisting}
The intrinsic input is built by:
\begin{lstlisting}[language=Lean]
def intrinsicInput :
    CoverIndexedZeroCompactRepresentedStokesIntrinsicInput
      (I := I) (omega := omega) supportCarrier where
  pointwise :=
    PointwiseCompactSupportChartBoxData.ofHalfSpaceModelSelfCharts
      (I := I) supportCarrier
  hK := X.compactSupport
  chartwiseSmooth := X.chartwiseSmooth
  support_subset_K := subset_closure
\end{lstlisting}
This is the first major compression: the four-field intrinsic theorem is
generated from the two-field first-principles theorem.

\subsection{Canonical represented endpoints}

The top-level endpoints are:
\begin{lstlisting}[language=Lean]
def canonicalRepresentedBulkIntegral : Real :=
  (X.intrinsicInput (I := I) (omega := omega)).
    canonicalRepresentedBulkIntegral

def canonicalRepresentedBoundaryIntegral : Real :=
  (X.intrinsicInput (I := I) (omega := omega)).
    canonicalRepresentedBoundaryIntegral
\end{lstlisting}
They inherit their meaning from the intrinsic route: finite selected cover,
selected smooth refinement, zero-localized local data, half-space local
Stokes, and endpoint reconstruction.

\subsection{The theorem}

The theorem is:
\begin{lstlisting}[language=Lean]
theorem representedStokes :
    X.canonicalRepresentedBulkIntegral (I := I) (omega := omega) =
      X.canonicalRepresentedBoundaryIntegral (I := I) (omega := omega) := by
  simpa [canonicalRepresentedBulkIntegral,
    canonicalRepresentedBoundaryIntegral] using
    (X.intrinsicInput (I := I) (omega := omega)).
      representedStokes
\end{lstlisting}
The proof is short because the intrinsic theorem has already generated the
route:
\begin{lstlisting}[language=Lean]
selectedSmoothRefinement
zeroSupportOfSelectedSmoothRefinement
interiorFieldsOfSelectedSmoothRefinement
hasIntrinsicRoute_intrinsic
representedStokes
\end{lstlisting}

\subsection{Intrinsic route}

The intrinsic structure still has four mathematical fields:
\begin{lstlisting}[language=Lean]
structure CoverIndexedZeroCompactRepresentedStokesIntrinsicInput
    (K : Set M) where
  pointwise : PointwiseCompactSupportChartBoxData I K
  hK : IsCompact K
  chartwiseSmooth : ManifoldForm.ChartwiseSmooth I omega
  support_subset_K : ManifoldForm.support I omega <= K
\end{lstlisting}
It is useful as an internal theorem-facing layer.  It states the theorem with a
chosen compact carrier and pointwise chart-box data.  The first-principles
theorem hides both by choosing \(K=\overline{\supp(\omega)}\) and generating
pointwise data from the half-space model.

The intrinsic route defines the selected cover, partition, and boundary finite
cover through:
\begin{lstlisting}[language=Lean]
openSelectedCover
selectedCover
openSelectedPartition
selectedPartition
selectedBoundaryFiniteCover
\end{lstlisting}
The canonical smooth refinement is generated by
\[
\code{selectedSmoothRefinement}.
\]
The two critical automatically generated proof obligations are:
\[
\code{zeroSupportOfSelectedSmoothRefinement}
\]
and
\[
\code{interiorFieldsOfSelectedSmoothRefinement}.
\]
Together they feed into
\[
\code{hasIntrinsicRoute_intrinsic}.
\]

\begin{theorem}[First-principles represented Stokes]
\label{thm:first-principles}
In the half-space model context, chartwise smoothness of \(\omega\) and
compactness of \(\overline{\supp(\omega)}\) imply equality of the canonical
represented bulk and boundary endpoints generated by the compact-support
local-to-global route.
\end{theorem}

The proof of \cref{thm:first-principles} is not an invocation of a black-box
global theorem.  It is the composition of the following generated statements:
compactness produces selected finite data; selected finite data produces a
smooth half-space refinement; the refinement produces zero-localized support
and analytic interior fields; local half-space Stokes gives each refined
identity; and endpoint reconstruction sums those identities.

\subsection{Reading the theorem mathematically}

Mathematically, the final theorem says:

\begin{quote}
For a chartwise smooth \(n\)-form on a half-space-modeled smooth manifold, if
the closure of its support is compact, then the finite coordinate bulk endpoint
generated by the compact-support Stokes proof equals the finite coordinate
boundary endpoint generated by the same proof.
\end{quote}

This is the clean statement the development was organized to reach.  The
large amount of internal code exists so that the theorem can be stated in this
small form without making the user provide the hidden proof objects.
